\documentclass[10pt,a4paper]{amsart}
\usepackage{amsmath,amssymb,amsthm,graphicx}
% For commutative diagrams we use Xy-pic
\usepackage[all]{xy}

% We use 2cell for 2-commutative diagrams.
\xyoption{2cell}
% Theorem environments.
%
\theoremstyle{plain}
\newtheorem{theorem}[subsection]{Theorem}
\newtheorem{proposition}[subsection]{Proposition}
\newtheorem{lemma}[subsection]{Lemma}

\theoremstyle{definition}
\newtheorem{definition}[subsection]{Definition}
\newtheorem{example}[subsection]{Example}
\newtheorem{exercise}[subsection]{Exercise}
\newtheorem{situation}[subsection]{Situation}

\theoremstyle{remark}
\newtheorem{remark}[subsection]{Remark}
\newtheorem{remarks}[subsection]{Remarks}

\numberwithin{equation}{subsection}

% Macros
%
\def\lim{\mathop{\mathrm{lim}}\nolimits}
\def\colim{\mathop{\mathrm{colim}}\nolimits}
\def\Spec{\mathop{\mathrm{Spec}}}
\def\Hom{\mathop{\mathrm{Hom}}\nolimits}
\def\Ext{\mathop{\mathrm{Ext}}\nolimits}
\def\SheafHom{\mathop{\mathcal{H}\!\mathit{om}}\nolimits}
\def\SheafExt{\mathop{\mathcal{E}\!\mathit{xt}}\nolimits}
\def\Sch{\mathit{Sch}}
\def\Mor{\mathop{\mathrm{Mor}}\nolimits}
\def\Ob{\mathop{\mathrm{Ob}}\nolimits}
\def\Sh{\mathop{\mathit{Sh}}\nolimits}
\def\NL{\mathop{N\!L}\nolimits}
\def\CH{\mathop{\mathrm{CH}}\nolimits}
\def\proetale{{pro\text{-}\acute{e}tale}}
\def\etale{{\acute{e}tale}}
\def\QCoh{\mathit{QCoh}}
\def\Ker{\mathop{\mathrm{Ker}}}
\def\Im{\mathop{\mathrm{Im}}}
\def\Coker{\mathop{\mathrm{Coker}}}
\def\Coim{\mathop{\mathrm{Coim}}}

% Boxtimes
%
\DeclareMathSymbol{\boxtimes}{\mathbin}{AMSa}{"02}

%
% Macros for moduli stacks/spaces
%
\def\QCohstack{\mathcal{QC}\!\mathit{oh}}
\def\Cohstack{\mathcal{C}\!\mathit{oh}}
\def\Spacesstack{\mathcal{S}\!\mathit{paces}}
\def\Quotfunctor{\mathrm{Quot}}
\def\Hilbfunctor{\mathrm{Hilb}}
\def\Curvesstack{\mathcal{C}\!\mathit{urves}}
\def\Polarizedstack{\mathcal{P}\!\mathit{olarized}}
\def\Complexesstack{\mathcal{C}\!\mathit{omplexes}}
% \Pic is the operator that assigns to X its picard group, usage \Pic(X)
% \Picardstack_{X/B} denotes the Picard stack of X over B
% \Picardfunctor_{X/B} denotes the Picard functor of X over B
\def\Pic{\mathop{\mathrm{Pic}}\nolimits}
\def\Picardstack{\mathcal{P}\!\mathit{ic}}
\def\Picardfunctor{\mathrm{Pic}}
\def\Deformationcategory{\mathcal{D}\!\mathit{ef}}
\newenvironment{reference}{}{}
\newenvironment{slogan}{}{}
\newenvironment{history}{}{}
\usepackage{cleveref}
\crefname{originallemma}{original lemma}{original lemmas}
\Crefname{originallemma}{Original lemma}{Original lemmas}
\makeatletter
\newlabel{lemma-divisor-special-fiber}{{0C5Z}{source}{}{sourcepointer-1}{}}
\newlabel{lemma-divisor-special-fiber@cref}{{[originallemma][0][]0C5Z}{[1][0][]source}{}{}{}}
\makeatother
\crefname{originallemma}{original lemma}{original lemmas}
\Crefname{originallemma}{Original lemma}{Original lemmas}
\makeatletter
\newlabel{lemma-intersection-pairing}{{0C64}{source}{}{sourcepointer-2}{}}
\newlabel{lemma-intersection-pairing@cref}{{[originallemma][0][]0C64}{[1][0][]source}{}{}{}}
\makeatother
\crefname{originalequation}{original equation}{original equations}
\Crefname{originalequation}{Original equation}{Original equations}
\makeatletter
\newlabel{equation-form}{{0C65}{source}{}{sourcepointer-3}{}}
\newlabel{equation-form@cref}{{[originalequation][0][]0C65}{[1][0][]source}{}{}{}}
\makeatother
\crefname{originallemma}{original lemma}{original lemmas}
\Crefname{originallemma}{Original lemma}{Original lemmas}
\makeatletter
\newlabel{divisors-lemma-ses-add-divisor}{{0C4T}{source}{}{sourcepointer-4}{}}
\newlabel{divisors-lemma-ses-add-divisor@cref}{{[originallemma][0][]0C4T}{[1][0][]source}{}{}{}}
\makeatother
\crefname{originallemma}{original lemma}{original lemmas}
\Crefname{originallemma}{Original lemma}{Original lemmas}
\makeatletter
\newlabel{divisors-lemma-sum-effective-Cartier-divisors-union}{{0C4R}{source}{}{sourcepointer-5}{}}
\newlabel{divisors-lemma-sum-effective-Cartier-divisors-union@cref}{{[originallemma][0][]0C4R}{[1][0][]source}{}{}{}}
\makeatother
\crefname{originalremark}{original remark}{original remarks}
\Crefname{originalremark}{Original remark}{Original remarks}
\makeatletter
\newlabel{divisors-remark-ses-regular-section}{{0C6K}{source}{}{sourcepointer-6}{}}
\newlabel{divisors-remark-ses-regular-section@cref}{{[originalremark][0][]0C6K}{[1][0][]source}{}{}{}}
\makeatother
\crefname{originallemma}{original lemma}{original lemmas}
\Crefname{originallemma}{Original lemma}{Original lemmas}
\makeatletter
\newlabel{morphisms-lemma-scheme-theoretic-union}{{0C4J}{source}{}{sourcepointer-7}{}}
\newlabel{morphisms-lemma-scheme-theoretic-union@cref}{{[originallemma][0][]0C4J}{[1][0][]source}{}{}{}}
\makeatother
\crefname{originallemma}{original lemma}{original lemmas}
\Crefname{originallemma}{Original lemma}{Original lemmas}
\makeatletter
\newlabel{varieties-lemma-euler-characteristic-additive}{{08AA}{source}{}{sourcepointer-8}{}}
\newlabel{varieties-lemma-euler-characteristic-additive@cref}{{[originallemma][0][]08AA}{[1][0][]source}{}{}{}}
\makeatother
\title{Adding a component to a divisor supported on the special fibre, and the resulting short exact sequence of cohomology}
\author{The Stacks Project Authors}
\date{stacks-project:models.tex:local-snapshot-20261012; GNU FDL 1.2 or later}
\begin{document}
\maketitle
\noindent\textit{Editorial note.} The counted unit is the Stacks Project Lemma lemma-add-component at source lines 5110-5121 together with its proof at source lines 5123-5154, a straight adjacent pairing in the chapter file. The chapter is the upstream file models.tex; the pool directory carries the longer name stacks-models and is not used as the chapter name anywhere in this card.\par\medskip
\noindent\textit{Editorial note.} One context span carries the regular-model situation at lines 4628-4637, which the retained statement names and which its own reference to the divisor lemma makes resolvable through a pointer; it is context rather than a counted unit and does not overlap the counted interval.\par\medskip
\noindent\textit{Editorial note.} The retained proof names three results of this chapter and four results of the Divisors, Morphisms and Varieties chapters. All seven are declared as non-load-bearing navigation pointers with their stable Stacks tags looked up in the shared tag index; the chapter ones are counted by cards of this lane and none of their intervals is reproduced here.\par\medskip
\noindent\textit{Editorial note.} The retained proof span starts at line 5123, the line of the opening marker read from the source, and contains exactly one proof environment.\par\medskip
\noindent\textit{Editorial note.} No citation occurs in the retained spans, so no bibliography entry is supplied.\par\medskip
\noindent\textit{Allowed standard prerequisites (editorial).}\begin{itemize}
\item Divisors supported on the special fibre, the intersection pairing and the additivity of the Euler characteristic. The situation fixes a discrete valuation ring, a smooth projective curve over its fraction field, a regular proper model and the components of the special fibre with their multiplicities; a divisor supported on the special fibre is a nonnegative combination of its components, the degree pairing of a combination of components against an invertible sheaf is symmetric, and for an effective Cartier divisor the structure sheaf of the sum of two summands fits into a short exact sequence comparing it with the sheaves of the summands and of their intersection. The Euler characteristic of a coherent sheaf on a proper scheme over a field is additive in short exact sequences, and the intersection numbers of the components are read off from the numerical type. The author's own notions and the quoted results of the Divisors, Morphisms and Varieties chapters are the allowed inputs.. The counted lemma is the bookkeeping step used in the genus formula of the section: it computes how the invariants change when one component is added to a divisor supported on the special fibre.
\end{itemize}\medskip
\ifdefined\hypertarget\hypertarget{sourcepointer-1}{}\fi
\noindent\textit{Source reference.} Original lemma 0C5Z: focused-seeds/AGP/sources/stacks-snapshot-20261011/tags.txt, source lines 15279--15279. Navigation pointer carried by the retained situation, which names this lemma when it writes the special fibre as a sum of its components with multiplicities; that lemma is counted by a separate card of this lane and its interval is not reproduced here. The number is the stable Stacks tag of the target, looked up in the shared tag index.\par\medskip
\ifdefined\hypertarget\hypertarget{sourcepointer-2}{}\fi
\noindent\textit{Source reference.} Original lemma 0C64: focused-seeds/AGP/sources/stacks-snapshot-20261011/tags.txt, source lines 15284--15284. Navigation pointer invoked by the retained proof for the symmetry of the degree pairing of the components, used to compute the shift in the intersection numbers; that lemma is counted by a separate card of this lane and its interval is not reproduced here. The number is the stable Stacks tag of the target, looked up in the shared tag index.\par\medskip
\ifdefined\hypertarget\hypertarget{sourcepointer-3}{}\fi
\noindent\textit{Source reference.} Original equation 0C65: focused-seeds/AGP/sources/stacks-snapshot-20261011/tags.txt, source lines 15285--15285. Navigation pointer to the chapter's displayed formula for the intersection numbers of the components, used by the retained proof; the cited labelled equation is an allowed input and its text is not reproduced here. The number is the stable Stacks tag of the target, looked up in the shared tag index.\par\medskip
\ifdefined\hypertarget\hypertarget{sourcepointer-4}{}\fi
\noindent\textit{Source reference.} Original lemma 0C4T: focused-seeds/AGP/sources/stacks-snapshot-20261011/tags.txt, source lines 15238--15238. Allowed input of the Divisors chapter quoted by the retained proof for the short exact sequence comparing an effective Cartier divisor with the sum of its summands; the pointer is registered so the reference resolves and the cited statement is not reproduced inside the span. The number is the stable Stacks tag of the target, looked up in the shared tag index.\par\medskip
\ifdefined\hypertarget\hypertarget{sourcepointer-5}{}\fi
\noindent\textit{Source reference.} Original lemma 0C4R: focused-seeds/AGP/sources/stacks-snapshot-20261011/tags.txt, source lines 15236--15236. Allowed input of the Divisors chapter quoted by the retained proof for the description of a sum of effective Cartier divisors as a union; the pointer is registered so the reference resolves and the cited statement is not reproduced inside the span. The number is the stable Stacks tag of the target, looked up in the shared tag index.\par\medskip
\ifdefined\hypertarget\hypertarget{sourcepointer-6}{}\fi
\noindent\textit{Source reference.} Original remark 0C6K: focused-seeds/AGP/sources/stacks-snapshot-20261011/tags.txt, source lines 15300--15300. Allowed input of the Divisors chapter quoted by the retained proof for the short exact sequence attached to a regular section of an invertible sheaf; the pointer is registered so the reference resolves and the cited remark is not reproduced inside the span. The number is the stable Stacks tag of the target, looked up in the shared tag index.\par\medskip
\ifdefined\hypertarget\hypertarget{sourcepointer-7}{}\fi
\noindent\textit{Source reference.} Original lemma 0C4J: focused-seeds/AGP/sources/stacks-snapshot-20261011/tags.txt, source lines 15229--15229. Allowed input of the Morphisms chapter quoted by the retained proof to identify the scheme-theoretic union of the two divisors; the pointer is registered so the reference resolves and the cited statement is not reproduced inside the span. The number is the stable Stacks tag of the target, looked up in the shared tag index.\par\medskip
\ifdefined\hypertarget\hypertarget{sourcepointer-8}{}\fi
\noindent\textit{Source reference.} Original lemma 08AA: focused-seeds/AGP/sources/stacks-snapshot-20261011/tags.txt, source lines 10525--10525. Allowed input of the Varieties chapter quoted by the retained proof for the additivity of the Euler characteristic in a short exact sequence, which is what turns the sequence into the numerical assertion; the pointer is registered so the reference resolves and the cited statement is not reproduced inside the span. The number is the stable Stacks tag of the target, looked up in the shared tag index.\par\medskip
\begin{situation}
\label{situation-regular-model}
Let $R$ be a discrete valuation ring with fraction field $K$,
residue field $k$, and uniformizer $\pi$.
Let $C$ be a smooth projective curve over $K$ with $H^0(C, \mathcal{O}_C) = K$.
Let $X$ be a regular proper model of $C$.
Let $C_1, \ldots, C_n$ be the irreducible components of the special
fibre $X_k$. Write $X_k = \sum m_i C_i$ as in
Lemma \ref{lemma-divisor-special-fiber}.
\end{situation}

\begin{lemma}
\label{lemma-add-component}
In Situation \ref{situation-regular-model} suppose we have an
effective Cartier divisors $D, D' \subset X$ such that
$D' = D + C_i$ for some $i \in \{1, \ldots, n\}$ and $D' \subset X_k$.
Then
$$
\chi(X_k, \mathcal{O}_{D'}) - \chi(X_k, \mathcal{O}_D) =
\chi(X_k, \mathcal{O}_X(-D)|_{C_i}) =
-(D \cdot C_i) + \chi(C_i, \mathcal{O}_{C_i})
$$
\end{lemma}

\begin{proof}
The second equality follows from the definition of the bilinear form
$(\ \cdot\ )$ in (\ref{equation-form}) and
Lemma \ref{lemma-intersection-pairing}. To see the first
equality we distinguish two cases.
Namely, if $C_i \not \subset D$, then $D'$ is the scheme
theoretic union of $D$ and $C_i$ (by
Divisors, Lemma \ref{divisors-lemma-sum-effective-Cartier-divisors-union})
and we get a short exact sequence
$$
0 \to \mathcal{O}_{D'} \to
\mathcal{O}_D \times \mathcal{O}_{C_i} \to
\mathcal{O}_{D \cap C_i} \to 0
$$
by Morphisms, Lemma \ref{morphisms-lemma-scheme-theoretic-union}.
Since we also have an exact sequence
$$
0 \to \mathcal{O}_X(-D)|_{C_i} \to
\mathcal{O}_{C_i} \to \mathcal{O}_{D \cap C_i} \to 0
$$
(Divisors, Remark \ref{divisors-remark-ses-regular-section})
we conclude that the claim holds
by additivity of euler characteristics
(Varieties, Lemma \ref{varieties-lemma-euler-characteristic-additive}).
On the other hand, if $C_i \subset D$ then we get an
exact sequence
$$
0 \to \mathcal{O}_X(-D)|_{C_i} \to \mathcal{O}_{D'} \to \mathcal{O}_D \to 0
$$
by Divisors, Lemma \ref{divisors-lemma-ses-add-divisor}
and we immediately see the lemma holds.
\end{proof}

\end{document}
